% =============================================================================
% Leakage-Aware Machine Learning Pipeline for Multi-Tenant ERP Systems
% IEEE-style two-column article. Compile with pdflatex + bibtex.
% =============================================================================
\documentclass[conference]{IEEEtran}

% --- Packages ----------------------------------------------------------------
\usepackage[utf8]{inputenc}
\usepackage[T1]{fontenc}
\usepackage{cite}
\usepackage{amsmath,amssymb,amsfonts}
\usepackage{algorithmic}
\usepackage{graphicx}
\usepackage{textcomp}
\usepackage{xcolor}
\usepackage{booktabs}
\usepackage{url}
\usepackage{listings}
\usepackage{hyperref}

\hypersetup{
    colorlinks=true,
    linkcolor=blue!50!black,
    citecolor=blue!50!black,
    urlcolor=blue!50!black,
}

\graphicspath{{figures/}}

% Listing style for code/SQL/Python snippets ----------------------------------
\lstset{
    basicstyle=\ttfamily\footnotesize,
    breaklines=true,
    columns=flexible,
    frame=single,
    framerule=0.4pt,
    xleftmargin=4pt,
    xrightmargin=4pt,
    aboveskip=4pt,
    belowskip=4pt,
    keywordstyle=\color{blue!60!black}\bfseries,
    commentstyle=\color{green!40!black}\itshape,
    stringstyle=\color{red!60!black},
    showstringspaces=false,
}

\begin{document}

% =============================================================================
\title{Leakage-Aware Machine Learning Integration in a\\Multi-Tenant ERP Platform: A Case-Study-Driven Approach}

\author{%
    \IEEEauthorblockN{[Author Name]}
    \IEEEauthorblockA{[Department / Institution]\\
    [City, Country]\\
    \texttt{[email@domain]}}
}

\maketitle

% =============================================================================
\begin{abstract}
Embedding machine learning (ML) inside enterprise resource planning (ERP)
systems is increasingly common, but published descriptions rarely address two
practical hazards: (i) target leakage that produces unrealistically optimistic
metrics, and (ii) per-tenant model lifecycle in shared-nothing multi-tenant
deployments. We describe an end-to-end ML pipeline integrated into a Django
multi-tenant ERP, covering CRM lead scoring, HR attrition prediction, accounts
receivable late-payment risk, and inventory demand forecasting. We treat
\emph{leakage detection} as a first-class validation gate using a label-shuffle
test, which rejected our initial lead-scoring features despite an apparently
strong area under the receiver operating characteristic curve (AUC) of
$0.9996$. The same gate validated the corrected pipeline. We further present a
detailed case study on late-payment prediction, where a naive
delivery-delay-based label was replaced with a customer-level payer profile;
this caused the headline AUC to fall from $0.63$ to $0.50$ before rising to
$0.78$ (5-fold cross-validation) once history-aware features became
identifiable through customer bucketing. We discuss what synthetic-data AUCs
do and do not measure, the design of a per-tenant model registry, and
limitations relevant to real-world deployment.
\end{abstract}

\begin{IEEEkeywords}
machine learning, ERP, multi-tenant SaaS, target leakage, demand forecasting,
late-payment prediction, customer churn, employee attrition, MLOps
\end{IEEEkeywords}

% =============================================================================
\section{Introduction}
\label{sec:intro}

Modern ERP platforms manage operational data covering customers, invoices,
products, employees, and inventory movements. Adding ML capabilities on top
of this data is a natural step: lead scoring for sales pipelines, attrition
risk for HR, late-payment risk for accounts receivable, and demand
forecasting for inventory replenishment are all standard supervised or
time-series problems with mature algorithms~\cite{lessmann2015benchmarking,
sakar2019machinelearning}.

In practice, however, integrating ML into a working multi-tenant ERP raises
problems that the algorithmic literature largely sidesteps:

\begin{itemize}
    \item \textbf{Target leakage from operational state.} ERP feature
    tables often contain post-event status fields (e.g., a lead's \texttt{won}
    flag, an invoice's \texttt{paid} timestamp). Naively including them in
    feature extraction yields models whose offline metrics are excellent and
    whose online performance collapses.
    \item \textbf{Per-tenant model lifecycle.} A shared-nothing multi-tenant
    SaaS deployment hosts each tenant's data in a separate database. A model
    trained on tenant $A$ should never score tenant $B$'s rows; the registry,
    feature pipeline, and serving layer must respect that boundary.
    \item \textbf{Sparse and irregular operational signals.} Most products
    sell rarely, most leads never convert, most employees do not leave, and
    most invoices are paid on time. Class imbalance and time-series sparsity
    are the rule, not the exception.
\end{itemize}

This paper reports an end-to-end design that addresses each of the issues
above, integrated into an open-source Django ERP we have developed. The
contributions are:

\begin{enumerate}
    \item A reproducible ML training pipeline integrated with a multi-tenant
    ORM and a per-tenant model registry (Section~\ref{sec:architecture}).
    \item A \emph{shuffle-label} validation gate used to detect and reject
    leaky feature definitions, with a worked example showing how it caught
    leakage hidden behind AUC $\approx 1.0$ (Section~\ref{sec:leakage}).
    \item Four case studies---lead scoring, attrition, late-payment risk, and
    demand forecasting---each documenting the data, label, features, leakage
    check, and limitations (Section~\ref{sec:cases}).
    \item A discussion of what synthetic-data evaluation does and does not
    establish, and how to report it honestly (Section~\ref{sec:discussion}).
\end{enumerate}

% =============================================================================
\section{Related Work}
\label{sec:related}

\textbf{ML for credit and payment risk.} Lessmann et al.~\cite{lessmann2015benchmarking}
benchmark $41$ classifiers across $8$ credit-scoring datasets, finding gradient
boosting and random forests competitive with deep models, and report typical
AUCs between $0.74$ and $0.88$ on real data. We use this range as a
plausibility band for our late-payment results.

\textbf{Employee attrition.} Sakar et al.~\cite{sakar2019machinelearning} and
follow-up work on the IBM HR dataset report AUCs in the $0.78$--$0.88$ range
when leakage from outcome-correlated fields (e.g., \texttt{OverTime} measured
\emph{after} the resignation decision) is controlled.

\textbf{Demand forecasting.} Prophet~\cite{taylor2018forecasting} provides a
practical decomposition baseline that handles weekly and yearly seasonality
without manual tuning, and is widely used as a benchmark in industrial
deployments.

\textbf{Target leakage.} Kaufman et al.~\cite{kaufman2012leakage} formalize
leakage in predictive modeling and propose audit techniques. The
\emph{label-shuffle} test we use is a special case of permutation testing
\cite{ojala2010permutation}: if a model trained on label-shuffled data still
exceeds chance AUC, the feature pipeline contains leakage independent of the
true label.

\textbf{Multi-tenant SaaS ML.} Industry write-ups discuss per-tenant versus
shared models~\cite{zhang2021multitenant}, but academic treatments remain
sparse. We adopt a per-tenant registry with shared training infrastructure.

% =============================================================================
\section{System Architecture}
\label{sec:architecture}

The host system is a Django~5 multi-tenant ERP with PostgreSQL backends.
Each tenant has a dedicated database, registered in a \texttt{TenantDatabase}
control table; a database router activates the correct connection within a
\texttt{TenantDatabaseContext} \cite{}. All ORM calls in feature-extraction
code use explicit \texttt{.using(db)} aliases, eliminating cross-tenant query
risk. Soft-deleted rows are accessible via an \texttt{all\_objects} manager
where ML history must include them (e.g., historical invoices for
late-payment features).

\subsection{Pipeline Stages}

\begin{figure}[t]
    \centering
    % \includegraphics[width=\columnwidth]{architecture.pdf}
    \fbox{\parbox{0.95\columnwidth}{\centering\textit{[Figure placeholder:
    architecture diagram --- tenant routing $\rightarrow$ feature extraction
    $\rightarrow$ leakage gate $\rightarrow$ training $\rightarrow$ per-tenant
    registry $\rightarrow$ scoring API.]}}}
    \caption{End-to-end pipeline from tenant-scoped feature extraction to
    per-tenant model serving.}
    \label{fig:architecture}
\end{figure}

\begin{enumerate}
    \item \textbf{Feature extraction.} Per-task selectors compose tenant-scoped
    Django querysets into \texttt{pandas} dataframes. Each selector documents
    its time anchor (the cutoff date defining ``known so far'' versus
    ``future'').
    \item \textbf{Leakage gate.} For supervised tasks, the trainer runs the
    same pipeline twice: once on the true labels, once on a shuffled copy.
    Hold-out AUC on shuffled labels must lie in $[0.45, 0.55]$; otherwise the
    run is rejected (Section~\ref{sec:leakage}).
    \item \textbf{Training.} Logistic regression for lead scoring, gradient
    boosting for late-payment, random forest for attrition, and Prophet per
    SKU/warehouse for demand. Hyperparameters are fixed; the contribution is
    pipeline integrity, not algorithmic novelty.
    \item \textbf{Persistence.} Models, metadata, and a training report are
    written under \texttt{ml\_models/\{company\_id\}/}, isolating artefacts
    per tenant.
    \item \textbf{Serving.} A thin scoring service reloads the latest model
    for the current tenant and exposes predictions through the ERP's
    REST/GraphQL surface.
\end{enumerate}

% =============================================================================
\section{Leakage as a First-Class Validation Gate}
\label{sec:leakage}

\subsection{Motivating Failure}

Our initial lead-scoring model achieved hold-out AUC $\approx 0.9996$. Manual
inspection of the feature extractor revealed two leaks: the \texttt{stage}
field was sampled \emph{after} the conversion event, and the activity-window
aggregator counted touchpoints whose timestamp followed the prediction
anchor. Both fixes were obvious in retrospect; neither was caught by a
standard time-aware train/test split because the offending fields were
correlated with the cutoff in subtle ways.

\subsection{The Shuffle-Label Test}

We adopted the test described in Algorithm~\ref{alg:shuffle}.
Implementation lives in \texttt{erp/ai/services/base.py} and is exposed via a
\texttt{--shuffle-labels} CLI flag on the training command.

\begin{figure}[h]
\begin{lstlisting}[language=Python]
def shuffle_label_auc(X, y, model_factory,
                      n_splits=5, seed=0):
    rng = np.random.default_rng(seed)
    y_perm = rng.permutation(y)
    cv = StratifiedKFold(n_splits=n_splits,
                         shuffle=True,
                         random_state=seed)
    scores = []
    for tr, te in cv.split(X, y_perm):
        m = model_factory().fit(X[tr], y_perm[tr])
        p = m.predict_proba(X[te])[:, 1]
        scores.append(roc_auc_score(y_perm[te], p))
    return float(np.mean(scores))
\end{lstlisting}
\caption{Label-shuffle leakage probe. Mean AUC outside $[0.45, 0.55]$ flags
the pipeline as leaky regardless of the real-label score.}
\label{alg:shuffle}
\end{figure}

\subsection{Outcome Across Models}

Table~\ref{tab:shuffle} summarises the gate results after the corrections
described in Section~\ref{sec:cases}. All four models pass the test.

\begin{table}[h]
\centering
\caption{Leakage gate: real vs.\ shuffled-label AUC.}
\label{tab:shuffle}
\begin{tabular}{lrr}
\toprule
\textbf{Model} & \textbf{Real AUC} & \textbf{Shuffled AUC} \\
\midrule
Lead scoring   & $0.9996$ & $0.4866$ \\
Attrition      & $0.9996$ & $0.5335$ \\
Late payment   & $0.7847$ (CV) / $0.8037$ (HO) & $0.5020$ \\
Demand (Prophet) & MAPE-$0.69$ to $0.74$ & n/a \\
\bottomrule
\end{tabular}
\end{table}

The gap between real and shuffled AUC is the only metric we treat as evidence
of learned signal. The absolute AUC, especially on synthetic data, is
discussed critically in Section~\ref{sec:discussion}.

% =============================================================================
\section{Case Studies}
\label{sec:cases}

\subsection{Late-Payment Risk (Featured Case Study)}
\label{sec:latepayment}

\textbf{Dataset.} The Olist Brazilian E-Commerce dataset
\cite{olist2018dataset} provides $\sim$$99{,}000$ orders with delivery
timestamps and payment records, imported into the ERP as customers,
invoices, payments, and product rows.

\textbf{Iteration 1: delivery-delay proxy (AUC $0.63$).} Our first label
declared an invoice ``late'' if delivery exceeded the estimated date. Hold-out
AUC was $0.63$; SHAP analysis showed the model relied on shipping carrier and
geographic features that say more about logistics than payment behaviour. The
target was the wrong concept.

\textbf{Iteration 2: hidden payer profile (AUC $0.50$).} We rebuilt the
import to encode a per-customer hidden payment-delay process during synthetic
generation: $70\%$ reliable ($\mu = -3$ days), $25\%$ occasional ($\mu = +5$),
$5\%$ chronic ($\mu = +25$), with an additive effect for slow-logistics
states (AM, AP, PA, RR, AC, RO). The label was now derived from a true
behavioural signal. AUC dropped to $0.50$. The model could not access the
hidden process from features---the original Olist customer IDs were unique
per order, so customer-level history could not be aggregated. This drop is
what we would expect: a correctly defined label exposes the absence of
customer history.

\textbf{Iteration 3: customer bucketing (AUC $0.78$).} We hashed the
$\sim$$96$k unique IDs into $3{,}000$ logical accounts. Customer-aggregate
features (\texttt{cust\_payment\_delay\_mean}, \texttt{cust\_late\_ratio},
\texttt{cust\_total\_invoices}, \texttt{amount\_vs\_cust\_avg}) became
informative because most accounts now had multiple invoices.

\begin{table}[h]
\centering
\caption{Late-payment iteration history.}
\label{tab:latepayment}
\begin{tabular}{lrr}
\toprule
\textbf{Iteration} & \textbf{AUC (hold-out)} & \textbf{Shuffled AUC} \\
\midrule
1. Delivery proxy      & $0.63$  & not run \\
2. Payer-profile label & $0.50$  & $\sim$$0.50$ \\
3. + customer bucketing & $0.80$  & $0.50$ \\
\bottomrule
\end{tabular}
\end{table}

The intermediate drop is, in our view, the most pedagogically valuable result
in this paper. It illustrates that a higher AUC is not always a better model;
the right question is whether the label encodes the concept of interest.

\subsection{Lead Scoring}

\textbf{Data.} A synthetic CRM lead generator produces leads with stages,
contact attempts, and source channels. \textbf{Label.} Won within $30$ days
of creation. \textbf{Fix.} Feature window is now symmetric, ending exactly at
the prediction anchor; \texttt{stage} is removed; activity counts use only
events with timestamp $<$ anchor.

\textbf{Result.} Real AUC $0.9996$, shuffled $0.49$. The high real AUC
reflects deterministic generator rules; we do not interpret it as expected
production performance (Section~\ref{sec:discussion}).

\subsection{Employee Attrition}

\textbf{Data.} IBM HR Employee Attrition~\cite{ibmhr}.
\textbf{Label.} \texttt{Attrition} flag.
\textbf{Fix.} Per-employee feature snapshot taken from a $90$-day backfilled
state preceding the resignation event, eliminating reliance on
post-resignation aggregates.

\textbf{Result.} Real AUC $0.9996$, shuffled $0.53$. Same caveat as lead
scoring---the labelled fields in the IBM dataset are deterministically
predictive when not deliberately perturbed.

\subsection{Demand Forecasting}

\textbf{Data.} Synthetic stock-movement series, $365$ days $\times$ $50$
products $\times$ $2$ warehouses, with three demand profiles (fast, moderate,
slow), weekly and monthly seasonality, and weekend dips.

\textbf{Method.} Prophet~\cite{taylor2018forecasting} fitted per
\textit{(SKU, warehouse)} pair. We previously observed that only $50\%$ of
\texttt{ReorderRule} entries trained successfully because the seeder
populated only the first warehouse. After fixing the seeder to iterate over
all active warehouses, $100/100$ rules trained.

\textbf{Backtest score.} Cross-validation MAPE was uninformative on slow
movers (Prophet skipped MAPE for $y\to 0$); we report mean absolute scaled
error in the range $[0.69, 0.74]$ across SKUs as the backtest signal. A
representative forecast is shown in Fig.~\ref{fig:prophet}.

\begin{figure}[t]
    \centering
    \fbox{\parbox{0.95\columnwidth}{\centering\textit{[Figure placeholder:
    Prophet backtest plot, fast vs.\ slow SKU.]}}}
    \caption{Prophet cross-validation forecast for representative fast and
    slow SKUs.}
    \label{fig:prophet}
\end{figure}

% =============================================================================
\section{Discussion}
\label{sec:discussion}

\subsection{What synthetic-data AUC measures}

Two of our four classification models attain AUC $\approx 0.9996$ on
synthetic data. Crucially, both pass the shuffle-label gate
(Table~\ref{tab:shuffle}), confirming the absence of leakage. But the
absolute number remains a pipeline-correctness claim, not a production
performance claim. Synthetic generators encode deterministic feature-to-label
rules with little aleatoric noise, so a sufficiently expressive model
recovers them almost perfectly.

For external comparison, real-world benchmarks in the same task
families~\cite{lessmann2015benchmarking, sakar2019machinelearning} report AUCs
in $0.74$--$0.88$. We therefore present our $0.9996$ values as
\emph{ceiling validation}---the pipeline can attain its theoretical maximum
on a clean signal---and \emph{not} as forecasts for tenant deployment.

\subsection{Why the late-payment number is different}

Late payment is the only model whose label was constructed with a hidden
stochastic process and whose features were realistically incomplete (only
aggregable history was available). This is closer to a real customer-risk
problem, and the resulting AUC of $0.78$--$0.80$ falls within the credit-risk
band reported in the literature.

\subsection{Per-tenant model lifecycle}

The registry layout \texttt{ml\_models/\{company\_id\}/} ensures that
artefacts and metadata never cross tenants. We considered an alternative
hierarchical model (one shared backbone with a per-tenant calibration head),
which trades training cost for inference complexity. For ERP tenants of
heterogeneous size and domain, we did not find a compelling argument for
sharing parameters across tenants; the per-tenant approach is also easier to
reason about under data-deletion requirements.

\subsection{Operational findings}

Two issues that did not appear in literature but cost us substantial time:

\begin{itemize}
    \item \textbf{ORM cascade deletion at scale.} Clearing $99$k Olist rows
    via Django's \texttt{.delete()} took $\sim$$30$ minutes due to in-memory
    PK loading. Replacing it with raw SQL \texttt{DELETE} statements (in
    correct foreign-key order, with explicit \texttt{StockMovement} and
    \texttt{ReorderRule} cascade) brought the clear phase under $2$ seconds.
    \item \textbf{Sparse seed coverage.} A demand seeder that populated only
    one of two warehouses caused $50\%$ of forecasts to fail silently with
    ``insufficient demand signal''. Iterating over all active warehouses
    fixed this.
\end{itemize}

% =============================================================================
\section{Limitations and Future Work}
\label{sec:limits}

\textbf{Single tenant evaluated.} All numbers in this paper come from one
synthetic tenant. Cross-tenant generalisation is unverified.

\textbf{Synthetic ceiling.} Our two highest-AUC models are bounded above by
the realism of the generators, not by the pipeline. Real-tenant pilots are
required to establish production performance.

\textbf{No drift monitoring.} The serving layer reloads the latest model but
does not yet track distribution shift between training and inference inputs.

\textbf{No online evaluation.} A/B testing against a no-ML baseline inside a
real tenant remains future work.

\textbf{Class imbalance handling.} We used default class weights and report
AUC. Threshold-dependent metrics (precision/recall at operating points) are
left to per-tenant tuning.

% =============================================================================
\section{Conclusion}
\label{sec:conclusion}

Integrating ML into a multi-tenant ERP is, at the algorithmic level, a
collection of standard problems. The interesting parts are operational:
designing feature pipelines that respect tenant isolation, treating leakage
detection as a deployment gate rather than a debugging convenience, and
distinguishing pipeline-correctness AUC from production-performance AUC. We
have shown a working pipeline that does this for four representative tasks,
and a late-payment case study whose iteration history is, in our view, more
valuable than any single AUC number.

% =============================================================================
\section*{Reproducibility}

The system, training command, and seed scripts are available at
\href{https://github.com/[user]/[repo]}{\texttt{[repo URL]}} under tag
\texttt{thesis-snapshot-v1}. Training reports for the experiments described
here are stored in \texttt{thesis/results/}.

% =============================================================================
\bibliographystyle{IEEEtran}
\bibliography{references}

\end{document}
